\section{Introduction}
\label{sec:intro}
Consider an npm asset library, its dependent UI package, a Docker image
packaging that UI, and a second image built from the first. The UI consumes
local output; an isolated registry-based image builder needs its base
published. A global build-then-publish schedule fails, while serializing
every package makes local-artifact consumers wait for unnecessary uploads.

If the final image fails after its providers publish, retry should retain its
release and target version. The engine must distinguish which versions are
owed, which stages may run, and which obligations survive failure.

Release engineering spans more than compilation~\cite{adams}. Without an atomic
commit across registries, execution resembles a saga~\cite{sagas}:
steps commit separately. Our protocol uses forward recovery because published
artifacts may not be retractable.

Existing tools schedule stages, compute versions, and coordinate
publication~\cite{lerna,nx,changesets,releaseplease,semanticrelease}.
We study their composition with a ledger that retains failed consumers'
obligations after providers succeed. Consumer-based admission keeps propagated
intent pending according to the consumer's own release window. Separate
windows distinguish prerelease-train intent from fresh, undelivered work.

We present \tool{} and its release-intent language \lang{}, using pseudonyms
for anonymous review; campaign records retain exact artifact identities.
Git history and tags form the ledger; ordinary commands build and publish.
No additional persistent release-state database is required, apart from
storage used by registries, build tools, and applications.

A package is a versioned unit of software, specifications, or documentation,
using common intent, dependency, command, and completion interfaces.
Guarantees require conforming integrations. The graph gives provisioned CI
a stable entry point for affected tests and releases; package configuration
still changes, and new credentials or runners may require workflow changes.

The paper makes three contributions:
\begin{itemize}
\item An ecosystem-independent, saga-inspired release protocol combining
local-artifact and published-artifact dependencies with explicit release
intent and per-package completion records, including a description of fleet
composition across repository-owned histories
(Section~\ref{sec:system}).
\item A self-contained model with deterministic planning and eight recovery,
discharge, and scope/channel guarantees under stated assumptions
(Sections~\ref{sec:model}--\ref{sec:guarantees}). Counterexamples explain why
consumer windows and stable retry predicates matter.
\item An evaluation of a released binary using differential planning,
targeted vectors, injected failures, historical replay, and scale measurements
(Section~\ref{sec:eval}), with the model-to-implementation gap made explicit.
\end{itemize}

The mathematical contribution is the admission and discharge discipline,
not a new topological-sorting or Git-storage primitive.
